\documentclass[11pt,a4paper]{article}

\usepackage{iftex}
\ifPDFTeX
  \usepackage[T1]{fontenc}
  \usepackage[utf8]{inputenc}
  \usepackage{lmodern}
\else
  \usepackage{fontspec}  % XeTeX/LuaTeX (e.g. tectonic): Unicode fonts, so °, á, í render directly
\fi
\usepackage[margin=2.3cm]{geometry}
\usepackage{amsmath,amssymb}
\usepackage{booktabs}
\usepackage{tabularx}
\usepackage{graphicx}
\newcolumntype{L}[1]{>{\raggedright\arraybackslash}p{#1}}
\newcolumntype{Y}{>{\raggedright\arraybackslash}X}
\usepackage{enumitem}
\usepackage{xcolor}
\usepackage[round,authoryear]{natbib}
\usepackage[colorlinks=true,linkcolor=blue!50!black,citecolor=blue!50!black,urlcolor=blue!50!black]{hyperref}

\newcommand{\code}[1]{\texttt{\detokenize{#1}}}
\newcommand{\assumption}[1]{\textcolor{orange!70!black}{[assumption]}~#1}

\title{Crop Rotations under Drought in the Valley of Mexico\\[0.3em]
\large Satellite data, a crop model with water and nitrogen, and the Mexican rainfed context}
\author{Gradient Descent Farmers --- NASA Space Apps warm-up\\
\small Working draft, branch \code{isaac/exploratory}}
\date{\today}

\begin{document}
\maketitle

\IfFileExists{tables/smap.tex}{\newcommand{\smapR}{0.87}
\newcommand{\smapDays}{731}
\newcommand{\smapRotation}{Maize every year}
}{}

\begin{abstract}
We want to compare crop sequences (rotations) for rainfed farms in central Mexico, where the rainy season is short,
variable and interrupted by a mid-summer dry spell (the \emph{canícula}). We built a pipeline that pulls free NASA
data (POWER weather, SMAP soil moisture, GLDAS and IGBP-DIS soil properties) for any point, and drives the WOFOST~8.1
crop model with it, adding nitrogen, frost, rain-onset sowing, fertilizer practice from a Texcoco field trial, and
season-to-season carryover of soil water and nitrogen. At CIMMYT's El Batán station (Texcoco) over 2005--2025, the
simulated fertilized maize yields (\(\approx\)6--7~t/ha dry matter) are in the range of field trials at Texcoco
(4.8--6.5~t/ha). Nitrogen, not water, is the largest yield gap in most years; rotating maize with faba bean raises maize
yields mostly in dry years; and frost and late sowing shorten the season more than expected. This document records the
data, equations, assumptions and results so the work can be reproduced and challenged.
\end{abstract}

\tableofcontents

%======================================================================
\section{Goal and approach}
%======================================================================

\paragraph{Decision problem.} Which crop to grow each rainy season, and in which order over the years, on a rainfed
(\emph{temporal}) farm in the Mexican highlands. The answer depends on the weather of each year, which nobody knows
when sowing, so we evaluate rotations over many real past years and under drier scenarios.

\paragraph{Approach.}
\begin{enumerate}[nosep]
  \item \textbf{Data} (Section~\ref{sec:data}): for a point (latitude, longitude), pull daily weather, soil moisture and
        static soil properties from NASA sources, and check that they are read correctly and behave physically.
  \item \textbf{Crop model} (Section~\ref{sec:model}): simulate each crop day by day with WOFOST~8.1 through the Python
        package PCSE, at three production levels: potential, water-limited and water- and nitrogen-limited.
  \item \textbf{Mexican context} (Section~\ref{sec:mexico}): sowing when the rains start, local fertilizer practice,
        legume nitrogen fixation, frost, drought indicators and drier-climate scenarios.
  \item \textbf{Rotations} (Section~\ref{sec:experiments}): chain seasons so water and nitrogen left by one crop affect
        the next, and compare rotations over 2005--2025.
\end{enumerate}

%======================================================================
\section{Study site}\label{sec:site}
%======================================================================

The default site is CIMMYT's El Batán research station in Texcoco, Estado de México (19.53\,°N, 98.85\,°W,
\(\approx\)2{,}250~m), in the farmland northeast of Mexico City, where long-term rainfed maize trials are run. The
climate has a rainy season from May to October and a dry, cool season from November to April with occasional frost.
In NASA POWER data, May--October rain averaged 774~mm over 2005--2025 (minimum 478~mm in 2023, maximum 1{,}005~mm in
2025), and June rain averaged 145~mm (12~mm in 2023). The soil at the site is a loam (GLDAS), with field capacity
0.373~m\(^3\)/m\(^3\), wilting point 0.137~m\(^3\)/m\(^3\) (IGBP-DIS) and porosity 0.437 (GLDAS), i.e.\ 236~mm of
plant-available water in the top metre.

%======================================================================
\section{Data}\label{sec:data}
%======================================================================

\begin{table}[htbp]
\centering\footnotesize
\caption{Data sources. All are free; SMAP and IGBP-DIS need a NASA Earthdata login.}
\label{tab:sources}
\begin{tabularx}{\linewidth}{L{2.6cm}YL{1.8cm}L{1.3cm}L{2.5cm}}
\toprule
Source & What it provides & Resolution & Format & Script \\
\midrule
NASA POWER \citep{nasapower} & Daily weather: min/max temperature, rain (bias-corrected), solar radiation, humidity, wind & 0.5°\(\times\)0.625° & JSON (API) & \code{power.py} \\
SMAP L4 SPL4SMGP v8 \citep{spl4smgp} & 3-hourly soil moisture (0--5~cm, 0--100~cm, profile), temperature, fluxes; 45 variables & 9~km EASE-Grid~2.0 & HDF5 & \code{smap.py} \\
GLDAS soils \citep{rodell2004gldas} & Texture class, sand/silt/clay, porosity & 0.25° & netCDF4 & \code{gldas_soils.py} \\
IGBP-DIS \citep{igbpdis2000} & Field capacity, wilting point, available water, bulk density, soil C and N (0--100~cm) & 1/12° (\(\approx\)9~km) & ASCII grid & \code{igbp_soils.py} \\
\bottomrule
\end{tabularx}
\end{table}

\subsection{Reading SMAP efficiently}
Each SMAP granule is a \(\approx\)140~MB HDF5 file covering the globe for one 3-hour window. Every variable is stored as
one gzip-compressed chunk of 1624\(\times\)3856 cells (\(\approx\)4.3~MB compressed), so even reading one pixel downloads
the whole chunk. Instead, \code{smap.py} asks NASA's OPeNDAP server for only the pixel
(\code{/Geophysical_Data/<var>[row][col]}), 16 requests in parallel, falling back to streaming the HDF5 if OPeNDAP
fails. This cut one week of data from \(\approx\)25~min to \(\approx\)30~s. The pixel is the cell whose centre is
nearest to the point; because EASE-Grid~2.0 is cylindrical, latitude is constant along rows and longitude along
columns, so the row and column are found from two 1-D arrays:
\begin{equation}
  i^* = \arg\min_i \lvert \varphi_i - \varphi \rvert, \qquad j^* = \arg\min_j \lvert \lambda_j - \lambda \rvert .
\end{equation}

\subsection{Validation of the data pipeline}
At Celaya, Guanajuato, for all of 2024 (2{,}928 three-hourly steps) we checked:
\begin{itemize}[nosep]
  \item OPeNDAP values equal direct HDF5 reads to 6 significant digits;
  \item the 1-D pixel lookup equals a full 2-D nearest search at 5 test sites (cell centres 2.8--4.8~km away);
  \item no missing or duplicate time steps; soil moisture within physical bounds;
  \item SMAP's precipitation input and POWER rain agree on the annual total (728 vs 727~mm), correlate daily
        (\(r = 0.69\)) and agree on rain/no-rain days 94\% of the time; surface soil moisture rises 0.019~m\(^3\)/m\(^3\)
        on days with \(>5\)~mm of POWER rain and falls on dry days;
  \item monthly SMAP land-surface temperature and POWER air temperature correlate at \(r = 0.99\).
\end{itemize}
SMAP L4 meets its accuracy requirement of 0.04~m\(^3\)/m\(^3\) unbiased RMSE against in-situ sensors
\citep{reichle2017smap}. Root-zone percentiles (\code{sm_rootzone_pctl}) are missing in NASA's files for about a third
of 2024 at Celaya, in dry-season blocks. SMAP data from 14 May to 28 July 2026 are being reprocessed after a
geolocation error \citep{nsidc2026notice}.

%======================================================================
\section{Crop model}\label{sec:model}
%======================================================================

We use WOFOST~8.1 \citep{dewit2019wofost,supit1994wofost} through PCSE~6.0.13 \citep{pcse}, with crop parameters from
the WOFOST crop parameter repository \citep{wofostcropparams}. WOFOST simulates one crop on one field, one day at a
time. Three model configurations give the three production levels used for yield-gap analysis:
\begin{center}
\begin{tabular}{lll}
\toprule
Level & PCSE model & Limiting factors \\
\midrule
Potential \(Y_p\) & \code{Wofost81_PP} & temperature, radiation, CO\(_2\) \\
Water-limited \(Y_w\) & \code{Wofost81_WLP_CWB} & + soil water (classic water balance) \\
Water- and N-limited \(Y_{wn}\) & \code{Wofost81_NWLP_CWB_CNB} & + nitrogen (classic N balance) \\
\bottomrule
\end{tabular}
\end{center}
Below, \(t\) is the day, \(T\) the mean daily temperature, and names in \code{TYPEWRITER} are WOFOST parameters.
A parameter ending in \code{TB} is a table \(\{(x_k, y_k)\}\) evaluated by linear interpolation.

\subsection{Weather: reference evapotranspiration}
PCSE computes the FAO-56 Penman--Monteith reference evapotranspiration \citep{allen1998fao56} from POWER data:
\begin{equation}
  ET_0 = \frac{\Delta\,R_n + \gamma\,\dfrac{900}{T+273}\,u_2\,(e_s - e_a)}{\Delta + \gamma\,(1 + 0.34\,u_2)},
\end{equation}
with \(R_n\) the net radiation (as mm of water), \(\Delta\) the slope of the saturation vapour pressure curve,
\(\gamma\) the psychrometric constant (which depends on air pressure, hence elevation), \(u_2\) wind at 2~m, and
\(e_s - e_a\) the vapour pressure deficit.

\subsection{Phenology}
The development stage \(D\) runs from 0 (emergence) through 1 (flowering) to 2 (maturity) and depends on temperature only:
\begin{equation}
  \frac{dD}{dt} =
  \begin{cases}
    \code{DTSMTB}(T)/\code{TSUM1}, & D < 1,\\
    \code{DTSMTB}(T)/\code{TSUM2}, & 1 \le D < 2.
  \end{cases}
\end{equation}
For the maize variety used, \code{DTSMTB} rises linearly from 0 at 6\,°C to 24 at 30\,°C, \code{TSUM1}=695 and
\code{TSUM2}=860\,°C\,d.

\subsection{Photosynthesis}
Leaf photosynthesis saturates with absorbed light \(I\):
\begin{equation}
  A(I) = A_{\max}\left(1 - e^{-\varepsilon I / A_{\max}}\right),
\end{equation}
with \(\varepsilon = \code{EFF}\). In WOFOST~8.1 the maximum rate depends on leaf nitrogen. The specific leaf N at
cumulative leaf area \(L\) from the top of the canopy is
\begin{equation}
  \mathrm{SLN}(L) = N_{LV}\,\frac{k_N\, e^{-k_N L}}{1 - e^{-k_N\,\mathrm{LAI}}},
  \qquad
  A_{\max}(L) = f_{CO_2}\, f_T\,\min\!\big(A_{\mathrm{ref}},\ \max(0,\ s\,(\mathrm{SLN}(L) - N_b))\big),
\end{equation}
with \(N_{LV}\) the leaf N (kg/ha), \(k_N=\code{KN}\), \(A_{\mathrm{ref}}=\code{AMAX_REF}\), \(s=\code{AMAX_SLP}\),
\(N_b=\code{AMAX_LNB}\). Light extinction through the canopy follows \(e^{-k L}\) for diffuse and direct light,
with sunlit and shaded leaves treated separately \citep{spitters1986daily}. Canopy photosynthesis is integrated with
3-point Gaussian quadrature over canopy depth and over the day:
\begin{equation}
  G = \mathrm{DAYL} \sum_{i=1}^{3} w_i\, \mathrm{LAI}\sum_{j=1}^{3} w_j\, A_{ij},
  \qquad x = (0.113,\ 0.5,\ 0.887),\ w = (0.278,\ 0.444,\ 0.278),
\end{equation}
where \(A_{ij}\) is the leaf rate at time point \(i\) and depth \(L = x_j\,\mathrm{LAI}\). The daily potential
assimilation in sugar is \(P = G\cdot\code{TMNFTB}(\bar T_{\min,7d})\cdot 30/44\), the last factor converting CO\(_2\) to CH\(_2\)O.

\subsection{Water stress and the soil water balance}
Maximum transpiration and the water-stress factor are
\begin{gather}
  T_{\max} = ET_0\left(1 - e^{-0.75\,k\,\mathrm{LAI}}\right),\\
  p = \frac{1}{0.76 + 1.5\,ET_0} - (5 - \code{DEPNR})\cdot 0.1,\qquad
  \theta_{cr} = (1-p)(\theta_{fc} - \theta_{wp}) + \theta_{wp},\\
  f_w = \operatorname{clip}\!\left(\frac{\theta - \theta_{wp}}{\theta_{cr} - \theta_{wp}},\ 0,\ 1\right),
\end{gather}
with \(\theta\) the root-zone moisture, \(\theta_{fc}\) field capacity and \(\theta_{wp}\) wilting point. Actual
transpiration is \(T_a = f_w T_{\max}\) and actual assimilation \(P\cdot f_w\). The classic water balance is a bucket
over the rooted depth \(R\) (growing \(\code{RRI}\)~cm/day up to 100~cm) plus a lower zone:
\begin{equation}
  \frac{dW}{dt} = \text{infiltration} - T_a - E_s - \text{percolation},\qquad \theta = W/R,
\end{equation}
where water above field capacity percolates to the lower zone and roots gain the lower zone's water as they deepen.
With no crop, \(R\) is the top 10~cm.

\subsection{Respiration, growth and partitioning}
\begin{gather}
  M = \big(r_R W_R + r_L W_L + r_S W_S + r_O W_O\big)\,\code{RFSETB}(D)\, Q_{10}^{(T - 25)/10},\\
  \Delta W_{\mathrm{tot}} = C\,(P f_w - M),\qquad
  C = \left[\left(\tfrac{F_L}{c_L} + \tfrac{F_S}{c_S} + \tfrac{F_O}{c_O}\right)(1 - F_R) + \tfrac{F_R}{c_R}\right]^{-1},
\end{gather}
with organ weights \(W\), maintenance coefficients \(r\), conversion efficiencies \(c\), and partitioning fractions
\(F_R, F_L, F_S, F_O\) from tables of \(D\). The grain grows by \(\Delta W_O = F_O (1 - F_R)\,\Delta W_{\mathrm{tot}}\),
and the yield is \(Y = \sum_t \Delta W_O\) (grain dry matter, \code{TWSO}).

\subsection{Nitrogen (WOFOST 8.1, classic N balance)}
Each organ demands N to reach its maximum concentration \(n^{\max}\):
\begin{equation}
  D_N = \sum_{\text{organs}} \Big[\max\big(n^{\max} W - N,\ 0\big) + n^{\max}\,\Delta W\Big].
\end{equation}
Legumes fix a fraction of the demand, and the rest is taken from the soil's mineral N pool \(N_{av}\), but only if
the crop is not severely water-stressed:
\begin{equation}
  F_N = \code{NFIX_FR}\cdot D_N,\qquad
  U_N = \begin{cases} \min\!\big(D_N - F_N,\ N_{av},\ \code{RNUPTAKEMAX}\big), & f_w > 0.01,\\ 0, & \text{otherwise.}\end{cases}
\end{equation}
Faba bean has \code{NFIX_FR}=0.8; maize and barley 0. The soil pool changes as
\begin{equation}
  \frac{dN_{av}}{dt} = \eta\,N_{\mathrm{fert}} + \text{mineralization} - U_N,
\end{equation}
with fertilizer recovery \(\eta\). Low leaf N lowers \(A_{\max}\), slows leaf expansion and speeds leaf death.

%======================================================================
\section{Mexican context added to the model}\label{sec:mexico}
%======================================================================

\subsection{Rainfed sowing}
Crops are sown when the rains start: on the first day from the crop's earliest sowing date \(t_0\) with at least
20~mm of rain over the previous five days, or after 45 days anyway:
\begin{equation}
  t_{\mathrm{sow}} = \min\Big\{t \ge t_0 : \sum_{k=0}^{4} P(t-k) \ge 20~\mathrm{mm}\Big\},\quad t_{\mathrm{sow}} \le t_0 + 45.
\end{equation}
A fixed-date alternative (\(t_{\mathrm{sow}} = t_0\)) is available.

\subsection{Crops and fertilizer}
\begin{table}[htbp]
\centering\footnotesize
\caption{Crops and management. Fertilizer recovery \(\eta = 0.5\) for all applications.}
\label{tab:crops}
\begin{tabularx}{\linewidth}{lL{2.9cm}L{1.3cm}L{1.6cm}L{1.1cm}Y}
\toprule
Crop & WOFOST variety & Earliest sowing & N (kg/ha) & Min.\ grain N & Basis \\
\midrule
Maize & \code{Grain_maize_202} & May 15 & 80 + 80 (day 0, 40) & 1.0\% & 160~kg N/ha: highest dose tested on rainfed hybrids at Texcoco \citep{alonso2023terra}; split \assumption{} \\
Barley & \code{Spring_barley_301} & July 1 & 60 (day 0) & 1.2\% & \assumption{no highland source found} \\
Faba bean & \code{Faba_bean_801} & June 15 & 40 (day 0) & 3.5\% & \assumption{} \\
\bottomrule
\end{tabularx}
\end{table}
WOFOST's spring-sown wheat is not available (its wheat parameters are winter wheats that need vernalization), and
common bean (\emph{frijol}), oats and intercropping (\emph{milpa}) are not modeled.

\subsection{Grain-N constraint}
WOFOST~8.1 does not bound the N concentration of the grain: in a test without fertilizer, maize made 14.5~t/ha of
grain holding 0.26\% N, which real maize does not do (typically 1.0--1.5\%). We therefore cap the N-limited yield by
the N actually allocated to the grain:
\begin{equation}
  Y_{wn} = \min\!\left(Y_{wn}^{\mathrm{WOFOST}},\ \frac{N_{\mathrm{grain}}}{n_{\min}}\right),
\end{equation}
with \(n_{\min}\) from Table~\ref{tab:crops}. Both values are kept in the outputs.

\subsection{Soil nitrogen supply}
Mineralization of soil organic N is taken from IGBP-DIS total N in 0--100~cm (12{,}709~kg N/ha at Texcoco):
\begin{equation}
  M_N = \mu\,N_{\mathrm{total}},\qquad \mu = 0.005~\text{yr}^{-1} \Rightarrow M_N \approx 64~\text{kg N/ha per season},
\end{equation}
\assumption{\(\mu\)}, released evenly over the first 150 days of each season (March--July).

\subsection{Chaining seasons: water and nitrogen carryover}
Each season (March~1 to end of February) is simulated separately, and its end state starts the next one:
\begin{align}
  \mathrm{WAV}_{s+1} &= W^{\mathrm{profile}}_{\mathrm{end},s} - R_{\max}\,\theta_{wp}, &
  \theta^{\mathrm{top}}_{s+1} &= \operatorname{clip}(\theta^{\mathrm{top}}_{\mathrm{end},s},\ \theta_{wp},\ \theta_{sat}),\\
  N_{av,s+1} &= \big(N_{av,\mathrm{end},s} + \rho\,N_{\mathrm{res},s}\big)(1 - \lambda), &
  N_{\mathrm{res}} &= N_{RT} + (W_{RT}^{\mathrm{dead}})\,n^{\mathrm{res}}_{RT} + (W_{LV}^{\mathrm{dead}})\,n^{\mathrm{res}}_{LV},
\end{align}
where the total profile water is conserved exactly across the boundary (checked numerically), \(N_{\mathrm{res}}\) is
the N in roots and dead leaves (stems and green leaves, the \emph{rastrojo}, are assumed removed for fodder, as is
common in Mexico), \(\rho = 0.5\) the share of residue N released the next season \assumption{}, and
\(\lambda = 0.5\) the share of carried mineral N lost before the next crop \assumption{}. Fertilizer recovery
\(\eta = 0.5\) follows the 30--50\% first-season recovery typical of cereals \citep{cassman2002nue}. The CO\(_2\)
concentration is \(380 + 2.2\,(y - 2005)\)~ppm.

\subsection{Frost}
The crop ends on the first night, from 30 days after sowing, with
\begin{equation}
  T_{\min}(t) + \delta \le 0\,^{\circ}\mathrm{C},
\end{equation}
where \(\delta\) (default 0) lets us represent valley floors colder than the 0.5° POWER grid (\(\delta = -2\)~°C in
one scenario). POWER shows on average 5.4 frost days per year at Texcoco (1990--2025), with first autumn frosts
between mid-October and late December and none at all in 2021--2024.

\subsection{Drought indicators and scenarios}
For each season we record: rain from sowing to harvest; July--August rain (\emph{canícula}); the longest run of days
with \(<1\)~mm of rain; water-stress days (\(f_w < 0.5\)); soil water on March~1. Years are classed as dry, normal or
wet by terciles of May--October rain over 2005--2025. Drier climates are explored by multiplying daily rain by a
factor, either all year or only in July--August.

%======================================================================
\section{Experiments}\label{sec:experiments}
%======================================================================

Four rotations, one crop per rainy season: maize every year; maize--barley; maize--faba bean;
maize--barley--faba bean. Every phase of each rotation is simulated, so each crop occurs every year
(8 sequences \(\times\) 21 seasons, 2005--2025, at three production levels). The first season is a spin-up and is
excluded. Scenarios: baseline; 20\% and 40\% less rain; 50\% less rain in July--August; frost check 2\,°C colder; no
fertilizer; fixed sowing on each crop's earliest date.

%======================================================================
\section{Results}\label{sec:results}
%======================================================================

\subsection{Baseline: where the yield goes}
\begin{table}[htbp]
\centering
\footnotesize
\setlength{\tabcolsep}{3.5pt}
\caption{Baseline results at Texcoco, 2006--2025 (20 seasons per crop and rotation). Yields are grain or seed dry matter in t/ha: $Y_{wn}$ water- and N-limited (with the grain-N constraint), $Y_w$ water-limited, $Y_p$ potential. CV: coefficient of variation of $Y_{wn}$ across years. Failed: seasons with $Y_{wn} < 1$ t/ha. Frost: seasons ended by frost before maturity. N uptake in kg N/ha.}
\label{tab:baseline}
\resizebox{\ifdim\width>\linewidth\linewidth\else\width\fi}{!}{%
\begin{tabular}{llrrrrrrrrr}
\toprule
Rotation & Crop & $Y_{wn}$ & CV (\%) & $Y_w$ & $Y_p$ & Water gap (\%) & N gap (\%) & Failed & Frost & N uptake \\
\midrule
Maize every year & maize & 6.3 & 42 & 10.8 & 11.7 & 8 & 42 & 1 & 12 & 144 \\
Maize--barley & barley & 6.1 & 5 & 11.7 & 11.8 & 1 & 48 & 0 & 2 & 98 \\
Maize--barley & maize & 6.4 & 33 & 11.6 & 11.7 & 1 & 45 & 0 & 12 & 147 \\
Maize--faba bean & faba bean & 9.7 & 6 & 9.9 & 10.0 & 1 & 2 & 0 & 5 & 88 \\
Maize--faba bean & maize & 6.9 & 34 & 11.5 & 11.7 & 2 & 40 & 0 & 12 & 154 \\
Maize--barley--faba bean & barley & 6.1 & 5 & 11.7 & 11.8 & 1 & 48 & 0 & 2 & 98 \\
Maize--barley--faba bean & faba bean & 9.8 & 4 & 10.0 & 10.0 & 0 & 2 & 0 & 5 & 87 \\
Maize--barley--faba bean & maize & 6.8 & 31 & 11.8 & 11.7 & 0 & 42 & 0 & 12 & 155 \\
\bottomrule
\end{tabular}}
\end{table}


Table~\ref{tab:baseline} summarizes the baseline. Fertilized rainfed maize yields 6.3--6.9~t/ha of grain dry matter
(\(\approx\)7.3--8.0~t/ha at 14\% moisture), the same order as the 4.8--6.5~t/ha measured for rainfed hybrids at
Texcoco \citep{alonso2023terra}. Water-limited maize would yield \(\approx\)11~t/ha and potential maize 11.7~t/ha, so
\textbf{nitrogen, not water, is the largest yield gap} (40--45\%); the water gap (0--8\%) is concentrated in a few
years. The grain-N constraint is active in 69\% of maize seasons (on average $-2.7$~t/ha relative to raw WOFOST) and
in 95\% of barley seasons, and never for faba bean, whose fixation keeps its seed N high. Frost ended 12 of the 20
maize seasons before maturity: with rain-onset sowing in late May or June, the long-cycle maize is still filling
grain when the first frosts arrive (late October to December). Barley and faba bean tonnages are not credible
(Section~\ref{sec:limitations}).

\subsection{Rotations}
\begin{table}[htbp]
\centering
\small

\caption{Maize yield ($Y_{wn}$, t/ha) by the crop grown the season before (baseline).}
\label{tab:prev}
\resizebox{\ifdim\width>\linewidth\linewidth\else\width\fi}{!}{%
\begin{tabular}{lrrrr}
\toprule
Previous crop & Seasons & Mean & Min & Max \\
\midrule
barley & 20 & 6.4 & 1.4 & 11.7 \\
faba bean & 40 & 6.9 & 1.4 & 12.7 \\
maize & 20 & 6.3 & 0.0 & 10.1 \\
\bottomrule
\end{tabular}}
\end{table}


Maize after faba bean yields 6.9~t/ha against 6.3~t/ha in monoculture (+0.6~t/ha, +10\%, Table~\ref{tab:prev}).
Paired by year, maize after faba bean is better by more than 0.5~t/ha in 11 of 20 years and worse in 4. The gain
comes from nitrogen: maize after faba bean starts the season with more mineral N (11.0 vs 7.5~kg/ha on average),
because faba bean draws 80\% of its N from the air and leaves N-rich roots and dead leaves. In three of the four worse
years (2006, 2013, 2017), maize after faba bean started with wetter soil, grew a larger canopy and spread its N over
more leaves and stems, leaving less N for the grain; this response runs through the grain-N constraint and should be
treated as uncertain. In 2024, maize after maize benefited from 74~kg/ha of fertilizer N left unused after the
2023 failure.

The largest single difference is 2023, the driest season in the record (Table~\ref{tab:years}): maize after maize
failed (0.0~t/ha) while maize after faba bean yielded 12.7~t/ha. In 2022 the rains came late, maize was sown on
15~June and only matured on 5~January, drying the profile; on 1~March 2023 it held 146~mm against a wilting point
of 137~mm. Maize sown on 15~May then met a nearly rainless June (12~mm) and spent 117 days under water stress.

\subsection{Drought: amount versus timing}
\begin{table}[htbp]
\centering
\small

\caption{Mean maize yield ($Y_{wn}$, t/ha) by year class. Classes are terciles of May--October rain over 2005--2025; ranges in mm.}
\label{tab:yearclass}
\resizebox{\ifdim\width>\linewidth\linewidth\else\width\fi}{!}{%
\begin{tabular}{lrrr}
\toprule
Rotation & Dry (478--747 mm, 7 yr) & Normal (754--814 mm, 7 yr) & Wet (846--1005 mm, 7 yr) \\
\midrule
Maize every year & 5.2 & 5.5 & 8.0 \\
Maize--barley & 7.1 & 5.1 & 7.0 \\
Maize--faba bean & 7.8 & 5.4 & 7.7 \\
Maize--barley--faba bean & 7.5 & 5.4 & 7.7 \\
\bottomrule
\end{tabular}}
\end{table}

\begin{table}[htbp]
\centering
\footnotesize
\setlength{\tabcolsep}{3.5pt}
\caption{Year by year, maize every year (baseline): rain from sowing to harvest (mm), July--August rain (mm), longest run of days with $<1$ mm rain during the season, days with water-stress factor $<0.5$, soil water in the 0--100 cm profile on March 1 (mm; wilting point = 137 mm), and maize yield ($Y_{wn}$, t/ha) in monoculture vs after faba bean.}
\label{tab:years}
\resizebox{\ifdim\width>\linewidth\linewidth\else\width\fi}{!}{%
\begin{tabular}{rlrrrrrrrrl}
\toprule
Year & Class & Sowing & Season rain & Jul--Aug rain & Dry spell & Stress days & Soil water Mar 1 & Maize (mono) & Maize after faba & Frost end \\
\midrule
2006 & normal & May 19 & 746 & 313 & 6 & 15 & 195 & 10.1 & 8.1 & yes \\
2007 & wet & May 26 & 796 & 392 & 6 & 13 & 365 & 7.2 & 8.0 & yes \\
2008 & normal & May 19 & 747 & 310 & 17 & 13 & 355 & 7.2 & 7.3 & yes \\
2009 & normal & May 16 & 827 & 219 & 9 & 12 & 289 & 6.8 & 7.4 &  \\
2010 & normal & Jun 12 & 736 & 453 & 13 & 13 & 362 & 1.7 & 1.7 & yes \\
2011 & normal & Jun 21 & 692 & 378 & 5 & 10 & 274 & 1.4 & 1.4 & yes \\
2012 & normal & Jun 04 & 744 & 406 & 26 & 12 & 343 & 4.7 & 4.8 & yes \\
2013 & wet & May 15 & 992 & 318 & 6 & 18 & 193 & 10.1 & 8.0 & yes \\
2014 & wet & May 15 & 908 & 269 & 9 & 12 & 358 & 6.9 & 7.8 & yes \\
2015 & dry & May 15 & 699 & 213 & 7 & 10 & 302 & 6.6 & 7.3 &  \\
2016 & normal & Jun 03 & 779 & 328 & 17 & 10 & 319 & 6.5 & 7.1 &  \\
2017 & wet & May 29 & 820 & 363 & 10 & 12 & 227 & 9.7 & 7.6 & yes \\
2018 & dry & Jun 10 & 732 & 263 & 12 & 10 & 281 & 6.4 & 7.2 & yes \\
2019 & dry & Jun 03 & 684 & 261 & 16 & 9 & 281 & 6.4 & 7.0 &  \\
2020 & dry & Jun 11 & 558 & 280 & 16 & 13 & 299 & 5.1 & 5.2 & yes \\
2021 & wet & May 15 & 802 & 286 & 15 & 9 & 273 & 7.0 & 7.6 &  \\
2022 & dry & Jun 15 & 535 & 246 & 22 & 17 & 252 & 6.8 & 7.5 &  \\
2023 & dry & May 15 & 424 & 266 & 28 & 117 & 146 & 0.0 & 12.7 &  \\
2024 & wet & Jun 20 & 864 & 383 & 16 & 10 & 200 & 7.9 & 7.0 &  \\
2025 & wet & May 25 & 963 & 296 & 13 & 10 & 199 & 7.1 & 7.7 & yes \\
\bottomrule
\end{tabular}}
\end{table}


Total season rain is a weak predictor of maize yield here (Table~\ref{tab:yearclass}). Monoculture maize did only
slightly worse in dry years (5.2~t/ha) than in normal years (5.5~t/ha), and rotations did \emph{better} in dry years
than in normal ones. Two reasons appear in Table~\ref{tab:years}: frost ended 71\% of normal and wet seasons but only
33\% of dry ones, and what hurts is \emph{when} rain fails rather than how much falls. 2023 (478~mm of May--October rain, a
28-day dry spell, 12~mm in June) failed after maize; 2022 (565~mm) did not. Rotations protected maize in dry years (7.1--7.8 vs
5.2~t/ha), when the soil water and nitrogen left by the previous crop matter most.

\subsection{Scenarios}
\begin{table}[htbp]
\centering
\small

\caption{Mean maize yield ($Y_{wn}$, t/ha) by scenario and rotation, 2006--2025; in parentheses, the number of failed maize seasons ($<1$ t/ha).}
\label{tab:scenarios}
\resizebox{\ifdim\width>\linewidth\linewidth\else\width\fi}{!}{%
\begin{tabular}{lrrrr}
\toprule
Scenario & Maize every year & Maize--barley & Maize--faba bean & Maize--barley--faba bean \\
\midrule
Baseline & 6.3 (1) & 6.4 (0) & 6.9 (0) & 6.8 (0) \\
20\% less rain & 5.3 (1) & 6.0 (1) & 6.4 (0) & 6.6 (0) \\
40\% less rain & 2.4 (10) & 2.8 (8) & 2.6 (9) & 3.0 (9) \\
Canícula: $-$50\% rain Jul--Aug & 5.8 (2) & 6.0 (1) & 6.5 (1) & 6.5 (0) \\
Frost check 2\,°C colder & 5.7 (1) & 5.9 (1) & 6.5 (1) & 6.5 (1) \\
No fertilizer & 1.9 (1) & 1.9 (0) & 2.4 (0) & 2.3 (0) \\
Fixed sowing (May 15) & 8.5 (0) & 8.3 (0) & 9.0 (0) & 9.0 (0) \\
\bottomrule
\end{tabular}}
\end{table}


\begin{itemize}[nosep]
  \item \textbf{20\% less rain}: maize monoculture loses 1.0~t/ha; rotations lose 0.3--0.5~t/ha. Rotations buffer
        moderate drying.
  \item \textbf{40\% less rain}: 8--10 of 20 maize seasons fail in every rotation; rotation cannot compensate.
  \item \textbf{Stronger canícula} (half the July--August rain): $-$0.3 to $-$0.5~t/ha. At this site the soil
        stores 236~mm of usable water and buffers a weaker mid-summer, unlike shallow soils further north.
  \item \textbf{Colder nights} (frost check $-2$\,°C): frost ends 85\% of maize seasons; $-$0.4 to $-$0.6~t/ha.
  \item \textbf{No fertilizer}: maize falls to 1.9--2.4~t/ha, faba bean is unaffected. The Texcoco trial found only
        a small response (5.2~t/ha without fertilizer vs 5.6 with 160~kg N/ha; \citealp{alonso2023terra}), so the
        model's soil N supply (\(\mu\)) is probably too low for that site. Calibrating it is the first priority.
  \item \textbf{Fixed sowing on May 15}: maize rises to 8.3--9.0~t/ha. Earlier sowing lengthens the season and
        escapes some frost. WOFOST does not represent failed emergence in a dry seedbed, so this gain is optimistic.
\end{itemize}

\subsection{Soil moisture against SMAP}
\IfFileExists{tables/smap.tex}{%
The simulated 0--100~cm soil moisture (\smapRotation, 2023--2024) correlates with SMAP root-zone moisture at
\(r = \smapR\) over \smapDays\ days.}{}
The day-to-day dynamics agree, but levels differ: in the 2023 rainy season the model's profile (0.13--0.19) was drier
than SMAP's (0.27--0.32), as expected for a field whose previous crop dried it, compared with a 9~km cell with its own
soil description.

%======================================================================
\section{Verification}\label{sec:verification}
%======================================================================

The core daily equations shared by WOFOST~7.2 and~8.1 (phenology, light interception and Gaussian integration,
water stress, maintenance respiration, conversion and partitioning, Penman--Monteith \(ET_0\)) were recomputed by
hand for one simulated day (7 August 2023, maize after barley at Texcoco, WOFOST~7.2) and matched PCSE to all
printed decimals: \(ET_0 = 5.108\)~mm, \(dD/dt = 0.013442\), \(G = 943.79\)~kg CO\(_2\)/ha, \(f_w = 0.916\),
\(T_a = 3.821\)~mm, \(M = 138.82\)~kg CH\(_2\)O/ha, \(\Delta W_O = 302.36\)~kg/ha. The season's grain yield equals the
sum of daily grain growth (17{,}700~kg/ha). The nitrogen-specific equations of WOFOST~8.1 have not yet been verified
by hand.

%======================================================================
\section{Limitations}\label{sec:limitations}
%======================================================================

\begin{itemize}[nosep]
  \item \textbf{Crop parameters are European.} Absolute yields of barley and especially faba bean are not credible
        (simulated faba bean seed \(\approx\)10~t/ha vs.\ 1--2~t/ha typical in Mexico); compare maize across rotations,
        not the partner crops' tonnages. Maize yields are in the range of a Texcoco trial \citep{alonso2023terra} but
        the model is not calibrated.
  \item \textbf{Several nitrogen parameters are assumptions} (\(\mu, \rho, \lambda, \eta, n_{\min}\), barley and faba
        bean doses); results on rotation benefits depend on them.
  \item \textbf{No pests, diseases, weeds, phosphorus, hail or lodging}; no tillage or residue-cover effects on
        evaporation.
  \item \textbf{Weather is coarse.} The POWER cell's elevation is 2{,}633~m vs.\ \(\approx\)2{,}250~m at the site;
        valley-floor frosts are under-represented.
  \item \textbf{Frost ends the crop on a single threshold}; WOFOST has no partial frost damage for these crops.
  \item \textbf{Soil hydraulic conductivity and percolation} use PCSE defaults, not site data.
  \item \textbf{One site.} Results do not transfer to other soils or climates without rerunning.
  \item The rain-scaling scenarios change only rain amounts, not temperature or radiation.
\end{itemize}

%======================================================================
\section{Next steps}\label{sec:next}
%======================================================================

\begin{enumerate}[nosep]
  \item Calibrate maize (and partner crops) against observed yields: INIFAP/CIMMYT trials at Texcoco, SIAP municipal
        statistics.
  \item Replace or bias-correct POWER temperatures with station data (SMN, Chapingo) for frost.
  \item Use SMAP root-zone moisture to initialize or correct soil water at sowing.
  \item Add supplemental irrigation as a drought adaptation option.
  \item Add prices and costs to score rotations, then frame rotation choice as a sequential decision problem
        for reinforcement learning (CropGym; \citealp{kallenberg2023cropgym}; cf.\ CyclesGym, \citealp{turchetta2022cyclesgym}).
\end{enumerate}

%======================================================================
\section{Reproducibility}\label{sec:repro}
%======================================================================

Code: \code{src/data_ingestion/} and \code{src/crop_models/wofost_rotations.py}. Python 3.13, PCSE 6.0.13.
From the repository root:
\begin{verbatim}
python -m src.data_ingestion.smap --lat 19.53 --lon -98.85 --start 2023-01-01 \
    --end 2024-12-31 --variables sm_rootzone sm_surface sm_rootzone_wetness \
    --out data/smap_texcoco_2023_2024.csv
M=src.crop_models.wofost_rotations; L=data/rotations
python -m $M --smap data/smap_texcoco_2023_2024.csv --out $L/baseline.csv
python -m $M --rain-factor 0.8 --out $L/rain80.csv
python -m $M --rain-factor 0.6 --out $L/rain60.csv
python -m $M --rain-factor 0.5 --rain-months 7 8 --out $L/canicula50.csv
python -m $M --frost-tmin-offset -2 --out $L/frost_minus2.csv
python -m $M --n-scale 0 --out $L/no_fertilizer.csv
python -m $M --sowing fixed --out $L/fixed_sowing.csv
python docs/technical_report/make_tables.py
\end{verbatim}

\bibliographystyle{plainnat}
\bibliography{references}

\end{document}
